\chapter{Stability and integral curvature estimates}\label{sec:stable-curvature-integrals}
Here $F:\Sigma^n\to\mathbb R^{n+1}$ is a smooth two-sided minimal immersion, $n\geq2$, and $u=|A|$. Stability means
\[
 \int_\Sigma u^2\zeta^2\leq\int_\Sigma|\nabla\zeta|^2
 \qquad(\zeta\in C_c^\infty(\Sigma)).
\]
If there is a boundary, all supports lie in the interior. This is the compact-support version of Theorem \ref{thm:second-variation}. The geometric input is stability together with the improved Kato inequality. We follow the Euclidean part of Schoen--Simon--Yau \cite{SSY75}, Sections 1--2, pp.~279--283, and supply the regularization at zero curvature explicitly. No pointwise curvature estimate or singular-varifold regularity theorem is assumed.

\begin{lemma}[Regularized Simons inequality]\label{lem:regularized-simons}
For $\varepsilon>0$ let $w_\varepsilon=(u^2+\varepsilon^2)^{1/2}$. On a smooth Euclidean minimal hypersurface,
\[
 w_\varepsilon\Delta w_\varepsilon+w_\varepsilon^4
 \geq\frac2n|\nabla w_\varepsilon|^2
\]
holds classically everywhere, including where $A=0$.
\end{lemma}
\begin{proof}
The norm identity gives
\[
 w_\varepsilon\Delta w_\varepsilon+u^4
 =|\nabla A|^2-|\nabla w_\varepsilon|^2.
\]
Where $u>0$, one has $\nabla w_\varepsilon=(u/w_\varepsilon)\nabla u$, and Lemma \ref{lem:cmc-improved-kato} implies
$|\nabla A|^2\geq(1+2/n)|\nabla w_\varepsilon|^2$. At a zero of $A$, the identity $\nabla(u^2)=2\langle A,\nabla A\rangle$ gives $\nabla w_\varepsilon=0$, so the same inequality holds. Finally $w_\varepsilon^4\geq u^4$. This is a smooth form of the Euclidean differential inequality underlying \cite{SSY75}, equation (1.34).
\end{proof}

\begin{proposition}[The weighted curvature estimate from stability]\label{prop:stable-weighted-curvature}
Assume the immersion is stable. For $0\leq q<\sqrt{2/n}$ there is a constant $C_{n,q}<\infty$ such that every smooth compactly supported real $\phi$ satisfies
\[
 \int_\Sigma\phi^2u^{2q+4}
 \leq C_{n,q}\int_\Sigma u^{2q+2}|\nabla\phi|^2.
\]
The constants can deteriorate as $q\uparrow\sqrt{2/n}$. Compare \cite{SSY75}, equations (1.36)--(1.38), in zero ambient curvature.
\end{proposition}
\begin{proof}
Write $w=w_\varepsilon$ and introduce finite nonnegative integrals
\[
 I=\int\phi^2w^{2q}|\nabla w|^2,\quad
 J=\int w^{2q+2}|\nabla\phi|^2,\quad
 L=\int\phi^2w^{2q+4},\quad
 E=\varepsilon^2\int\phi^2w^{2q+2},
\]
and the signed cross term $K=\int\phi w^{2q+1}\langle\nabla w,\nabla\phi\rangle$. Cauchy--Schwarz gives $|K|\leq\sqrt{IJ}$. Multiply Lemma \ref{lem:regularized-simons} by $\phi^2w^{2q}$ and integrate by parts to obtain
\[
 (1+2q+2/n)I\leq L-2K.
\]
Use the smooth compactly supported test $\zeta=\phi w^{q+1}$ in stability. Since $u^2=w^2-\varepsilon^2$, this gives
\[
 L\leq(q+1)^2I+J+2(q+1)K+E.
\]
Combining the two inequalities, with $\delta=2/n-q^2>0$, yields
\[
 \delta I\leq J+2qK+E
 \leq J+\frac\delta2 I+\frac{2q^2}{\delta}J+E.
\]
Thus $I\leq(2/\delta+4q^2/\delta^2)J+(2/\delta)E$. A second Young inequality in the stability bound gives $L\leq2(q+1)^2I+2J+E$, hence $L\leq C_{n,q}(J+E)$ after increasing the constant. On the fixed compact support, $u$ is bounded, $E\to0$, and the other integrands converge by dominated convergence as $\varepsilon\downarrow0$. The resulting limit is the statement. This argument never differentiates $u$ at a zero without regularization.
\end{proof}

\begin{theorem}[The Schoen--Simon--Yau integral estimate in Euclidean space]\label{thm:ssy-integral-curvature}
For a stable smooth two-sided Euclidean minimal hypersurface of dimension $n\geq2$ and
\[
 4\leq p<4+\sqrt{8/n},
\]
there is $C_{n,p}<\infty$ such that every nonnegative compactly supported locally Lipschitz function $f$ in the interior satisfies
\[
 \int_\Sigma u^p f^p\leq C_{n,p}\int_\Sigma|\nabla f|^p.
\]
This is the Euclidean specialization of \cite{SSY75}, Theorem 1, pp.~281--282.
\end{theorem}
\begin{proof}
Put $q=(p-4)/2$, so the exponent assumption is exactly $0\leq q<\sqrt{2/n}$. First take smooth nonnegative $f$ and apply Proposition \ref{prop:stable-weighted-curvature} with $\phi=f^{p/2}$. If the power is not smooth at zeros, it is still compactly supported and continuously differentiable; smooth approximation in $W^{1,2}$ on a fixed compact neighborhood justifies the weighted inequality because $u$ is bounded there. Write $L=\int u^pf^p$. The chain rule and Holder give
\[
 L\leq C_{n,q}\frac{p^2}{4}\int(uf)^{p-2}|\nabla f|^2
 \leq C_{n,q}\frac{p^2}{4}L^{(p-2)/p}
             \left(\int|\nabla f|^p\right)^{2/p}.
\]
If $L=0$, there is nothing to prove. Otherwise divide by its indicated power and raise to $p/2$, absorbing the coefficient into $C_{n,p}$. For a compactly supported locally Lipschitz $f$, smooth nonnegative approximation in $W^{1,p}$ on a fixed compact neighborhood gives the statement by passage to the limit. Such approximation follows from coordinate mollification and a smooth partition of unity; the potential $u$ stays bounded on that neighborhood. No area or volume bound is used in this estimate.
\end{proof}

\begin{corollary}[Curvature integrals on intrinsic or extrinsic balls]\label{cor:stable-curvature-ball-integral}
Let $r:\Sigma\to[0,\infty)$ be a locally Lipschitz exhaustion with $|\nabla r|\leq1$ almost everywhere and compact sublevel sets, and suppose $\Sigma$ has no boundary. For $R>0$, $0<\theta<1$, and $4\leq p<4+\sqrt{8/n}$,
\[
 \int_{\{r<\theta R\}}u^p
 \leq\frac{C_{n,p}}{(1-\theta)^pR^p}\operatorname{Vol}(\{r<R\}).
\]
One may use intrinsic distance from a point on a complete connected hypersurface. One may instead use $r=|F-a|$ when $F$ is proper. Volumes refer to the induced metric on the parameter manifold, so immersion multiplicities are counted. Compare \cite{SSY75}, equation (2.9), p.~283.
\end{corollary}
\begin{proof}
Take $f=1$ for $r\leq\theta R$, let $f$ decrease linearly to zero between $\theta R$ and $R$, and set $f=0$ for $r\geq R$. It is a compactly supported Lipschitz test, with $|\nabla f|\leq((1-\theta)R)^{-1}$ almost everywhere in the annulus and zero elsewhere. Theorem \ref{thm:ssy-integral-curvature} gives the estimate. Intrinsic distance is 1-Lipschitz, and its closed balls are compact by completeness and Hopf--Rinow. For ambient distance, tangential projection gives the gradient bound and properness gives compact sublevel sets. Completeness alone is not being used to justify extrinsic compact support.
\end{proof}

\begin{theorem}[Subcritical volume growth forces flatness]\label{thm:stable-subcritical-growth}
In the setting of Corollary \ref{cor:stable-curvature-ball-integral}, suppose for some $p\in[4,4+\sqrt{8/n})$ there is a sequence $R_j\to\infty$ with
\[
 R_j^{-p}\operatorname{Vol}(\{r<R_j\})\longrightarrow0.
\]
Then $A=0$ everywhere. If the immersion is complete and connected, its image is an affine hyperplane. This is the Euclidean growth implication in \cite{SSY75}, Theorem 2, p.~283; the stated sequence assumption suffices for the proof.
\end{theorem}
\begin{proof}
Use $\theta=1/2$ in Corollary \ref{cor:stable-curvature-ball-integral}. Every fixed compact subset lies in $\{r<R_j/2\}$ for all sufficiently large $j$. Its integral of $u^p$ is therefore bounded by a quantity tending to zero. Continuity gives $u=0$ at every point. The chosen normal has zero derivative and hence is constant on a connected parameter manifold, so the image lies in an affine hyperplane. To see that it fills that hyperplane under completeness, start at one point and take a geodesic with any prescribed tangent vector. Its ambient acceleration is zero because both its intrinsic acceleration and the second fundamental form vanish. Completeness extends the geodesic to every time, giving every point on the corresponding straight line. The tangent space already equals the hyperplane's tangent space, so these lines cover the plane.
\end{proof}

\begin{corollary}[Euclidean-growth rigidity in dimensions at most five]\label{cor:stable-euclidean-growth-five}
Let $2\leq n\leq5$. A complete connected smooth two-sided stable minimal immersion in $\mathbb R^{n+1}$ has image an affine hyperplane if either its intrinsic balls have volume at most $CR^n$ for all $R\geq1$, or the immersion is proper and its extrinsic balls about a fixed center have induced volume at most $CR^n$ for all $R\geq1$.
\end{corollary}
\begin{proof}
Choose $p\in[4,4+\sqrt{8/n})$ with $p>n$. Such an exponent exists for $2\leq n\leq5$. The assumed growth gives $R^{-p}\operatorname{Vol}(\{r<R\})\leq CR^{n-p}\to0$, so Theorem \ref{thm:stable-subcritical-growth} applies with the appropriate exhaustion. For $n\geq6$, the interval has no exponent exceeding $n$, so this proof does not give the same conclusion from Euclidean growth. That limitation of the estimate is not a claim of a counterexample in dimension six. This is the dimension restriction explained after \cite{SSY75}, Theorem 2.
\end{proof}

\begin{corollary}[Entire minimal graphs through dimension five]\label{cor:bernstein-through-five}
For $2\leq n\leq5$, every $C^2$ solution $v:\mathbb R^n\to\mathbb R$ of the scalar minimal surface equation is affine. This is the graph consequence of the integral method in \cite{SSY75}, the remarks after Theorem 2; it is not the full dimension range of the Bernstein theorem.
\end{corollary}
\begin{proof}
The graph is smooth by Proposition \ref{prop:classical-graph-bootstrap}, two-sided, and stable by Corollary \ref{cor:graph-positive-jacobi}. It is proper since bounded ambient coordinates bound the horizontal coordinates and preimages of closed sets are closed. Its induced metric is complete: it dominates the Euclidean metric, and on any compact base set its smooth coefficients are bounded. Proposition \ref{prop:entire-graph-area-growth} gives Euclidean volume growth without a global gradient bound. Corollary \ref{cor:stable-euclidean-growth-five} makes its image an affine hyperplane. A hyperplane that is a graph over all of $\mathbb R^n$ is the graph of an affine function. The earlier two-dimensional proof in Theorem \ref{thm:bernstein-plane-graph} retains its separate logarithmic-cutoff argument.
\end{proof}
